\documentclass[UTF8,11pt]{ctexart}
\usepackage[a4paper,margin=22mm]{geometry}
\usepackage{graphicx,enumitem,amssymb,tabularx}
\usepackage[normalem]{ulem}
\usepackage{xeCJKfntef}
\setlength{\parindent}{0pt}
\setlength{\parskip}{.6em}
\sloppy
\begin{document}
\section*{2014年6月英语六级真题(第2套)}
\subsection*{2014 年6 月大学英语六级考试真题(二)}


\subsection*{\textbf{Part I} \textbf{Writing} \textbf{(30 minutes)}}


\textbf{Directions}：\textit{For this part},\textit{ you are allowed 30 minutes to write an essay explaining}\textbf{\textit{ why it is unwise to put all your eggs in one}} \textbf{\textit{basket}}.\textit{ You can give examples to illustrate your point}.\textit{ You should write at least}\textbf{\textit{ }}\textbf{\textit{\uline{150}}}\textit{ words but no more than 200 words.} 注意：此部分试题请在答题卡1 上作答。


\textbf{Part}II \textbf{Listening Comprehension} \textbf{(30 minutes)}


\subsection*{\textbf{Section A}}


\textbf{Directions}：\textit{In this section},\textit{ you will hear 8 short conversations and 2 long conversations. At the end of each conversation,} \textit{one or more questions will be asked about what was said. Both the conversation and the questions will be spoken only once}. \textit{After each question there will be a pause}.\textit{ During the pause, you must read the four choices marked A) 9 B) 9} C)\textit{ and} D),\textit{ and} \textit{decide which is the best answer. Then mark the}


\subsection*{\textit{corresponding letter on}\textbf{\textit{ Answer Sheet 1}}\textit{ with a single line through the centre} 注意：此部分试题请在答题卡1 上作答。}


1.


\begin{itemize}[leftmargin=2em,itemsep=.2em,topsep=.4em]\item[A.] They might be stolen goods.
\item[B.] They might be fake products.
\item[C.] They might be faulty products.
\item[D.] They might be smuggled goods.
\end{itemize}

2.


\begin{itemize}[leftmargin=2em,itemsep=.2em,topsep=.4em]\item[A.] They are civil servants.
\item[B.] They are job applicants.
\item[C.] They are news reporters.
\item[D.] They are public speakers.
\end{itemize}

3.


\begin{itemize}[leftmargin=2em,itemsep=.2em,topsep=.4em]\item[A.] The man has decided to quit his computer class.
\item[B.] The woman wants to get a degree in administration.
\item[C.] A computer degree is a must for administrative work.
\item[D.] The man went to change the time of his computer class.
\end{itemize}

4.


\begin{itemize}[leftmargin=2em,itemsep=.2em,topsep=.4em]\item[A.] A lot of contestants participated in the show.
\item[B.] The fifth contestant won the biggest prize.
\item[C.] It was not as exciting as he had expected.
\item[D.] It was sponsored by a car manufacturer.
\end{itemize}

5.


\begin{itemize}[leftmargin=2em,itemsep=.2em,topsep=.4em]\item[A.] Reading a newspaper column.
\item[B.] Looking at a railway timetable.
\item[C.] Driving from New York to Boston.
\item[D.] Waiting for someone at the airport.
\end{itemize}

6.


\begin{itemize}[leftmargin=2em,itemsep=.2em,topsep=.4em]\item[A.] He wears a coat bought in the mall.
\item[B.] He got a new job at the barbershop.
\item[C.] He had a finger hurt last night.
\item[D.] He had his hair cut yesterday.
\end{itemize}

7.


\begin{itemize}[leftmargin=2em,itemsep=.2em,topsep=.4em]\item[A.] He cannot appreciate the Picasso exhibition.
\item[B.] Even his nephew can draw as well as Picasso.
\item[C.] He is not quite impressed with modern paintings.
\item[D.] Some drawings by kindergarten kids are excellent.
\end{itemize}

8.


\begin{itemize}[leftmargin=2em,itemsep=.2em,topsep=.4em]\item[A.] He should not put the cart before the horse.
\item[B.] His conduct does not square with his words.
\item[C.] His attitude to student government has changed.
\item[D.] He has long been involved in student government.
\end{itemize}

\textbf{Questions 9 to 11 are based on the conversation you have just heard.}


9.


\begin{itemize}[leftmargin=2em,itemsep=.2em,topsep=.4em]\item[A.] She left her own car in Manchester.
\item[B.] Something went wrong with her car.
\item[C.] She wants to go traveling on the weekend.
\item[D.] Her car won’t be back in a week’s time.
\end{itemize}

\noindent\begin{tabularx}{\linewidth}{@{}lXXXX@{}}10. & \textbf{A.} Safety. & \textbf{B.} Comfort. & \textbf{C.} Size. & \textbf{D.} Cost\\\end{tabularx}\par

\noindent\begin{tabularx}{\linewidth}{@{}lXXXX@{}}11. & \textbf{A.} Third-party insurance. & \textbf{B.} Value-added tax & \textbf{C.} Petrol & \textbf{D.} CDW\\\end{tabularx}\par

\textbf{Questions 12 to 15 are based on the conversation you have just heard.}


12.


\begin{itemize}[leftmargin=2em,itemsep=.2em,topsep=.4em]\item[A.] How to update the basic facilities.
\item[B.] What to do to enhance their position.
\item[C.] Where to locate their plant.
\item[D.] How to attract investments.
\end{itemize}

13.


\begin{itemize}[leftmargin=2em,itemsep=.2em,topsep=.4em]\item[A.] Their road link to other European countries is fast.
\item[B.] They are all located in the south of France.
\item[C.] They are very close to each other.
\item[D.] Their basic facilities are good.
\end{itemize}

14.


\begin{itemize}[leftmargin=2em,itemsep=.2em,topsep=.4em]\item[A.] Try to avoid making a hasty decision.
\item[B.] Take advantage of the train links.
\item[C.] Talk with the local authorities.
\item[D.] Conduct field surveys first.
\end{itemize}

15.


\begin{itemize}[leftmargin=2em,itemsep=.2em,topsep=.4em]\item[A.] Future product distribution.
\item[B.] Local employment policies.
\item[C.] Road and rail links for small towns.
\item[D.] Skilled workforce in the hilly region.
\end{itemize}

\subsection*{\textbf{Section B}}


\textbf{Directions}：\textit{In this section},\textit{ you will hear 3 short passages. At the end of each passage},\textit{ you will hear some questions. Both} \textit{the passage and the questions will be spoken only once}.\textit{ After you hear a question},\textit{ you must choose the best answer from}


\subsection*{\textit{the four choices marked A), B) ,} C)\textit{ and D). Then mark the corresponding letter on}\textbf{\textit{ Answer Sheet 1}}\textit{ with a single line}}


\textit{through the centre.} 注意：此部分试题请在答题卡1 上作答。


\subsection*{\textbf{Passage One}}


\textbf{Questions 16 to 18 are based on the passage you have just heard.}


16.


\begin{itemize}[leftmargin=2em,itemsep=.2em,topsep=.4em]\item[A.] One fifth of them were on bad terms with their sisters and brothers.
\item[B.] About one eighth of them admitted to lingering bitter feelings.
\item[C.] More than half of them were involved in inheritance disputes.
\item[D.] Most of them had broken with their sisters and brothers.
\end{itemize}

17.


\begin{itemize}[leftmargin=2em,itemsep=.2em,topsep=.4em]\item[A.] Less concern with money matters.
\item[B.] More experience in worldly affairs.
\item[C.] Advance in age.
\item[D.] Freedom from work.
\end{itemize}

18.


\begin{itemize}[leftmargin=2em,itemsep=.2em,topsep=.4em]\item[A.] They have little time left to renew contact with their brothers and sisters.
\item[B.] They tend to forget past unhappy memories and focus on their present needs.
\item[C.] They are more tolerant of one another.
\item[D.] They find close relatives more reliable.
\end{itemize}

\subsection*{\textbf{Passage Two}}


\textbf{Questions 19 to 22 are based on the passage you have just heard.}


19.


\begin{itemize}[leftmargin=2em,itemsep=.2em,topsep=.4em]\item[A.] They have bright colors and intricate patterns.
\item[B.] They can only survive in parts of the Americas.
\item[C.] They are the only insect that migrates along fixed routes.
\item[D.] They have strong wings capable of flying long distances.
\end{itemize}

20.


\begin{itemize}[leftmargin=2em,itemsep=.2em,topsep=.4em]\item[A.] In a Michigan mountain forest.
\item[B.] In a Louisiana mountain forest.
\item[C.] In a Kentucky mountain forest.
\item[D.] In a Mexican mountain forest.
\end{itemize}

21.


\begin{itemize}[leftmargin=2em,itemsep=.2em,topsep=.4em]\item[A.] Each flock of butterflies lays eggs in the same states.
\item[B.] They start to lay eggs when they are nine months old.
\item[C.] Each generation in a cycle lays eggs at a different place.
\item[D.] Only the strongest can reach their destination to lay eggs.
\end{itemize}

22.


\begin{itemize}[leftmargin=2em,itemsep=.2em,topsep=.4em]\item[A.] Evolution of monarch butterflies.
\item[B.] Living habits of monarch butterflies.
\item[C.] Migration patterns of monarch butterflies.
\end{itemize}

\begin{itemize}[leftmargin=2em,itemsep=.2em,topsep=.4em]\item[D.] Environmental impacts on monarch butterfly life.
\end{itemize}

\subsection*{\textbf{Passage Three}}


\textbf{Questions 23 to 25 are based on the passage you have just heard.}


23.


\begin{itemize}[leftmargin=2em,itemsep=.2em,topsep=.4em]\item[A.] Time has become more limited.
\item[B.] Time has become more precious.
\item[C.] Time is money.
\item[D.] Time is relative.
\end{itemize}

24.


\begin{itemize}[leftmargin=2em,itemsep=.2em,topsep=.4em]\item[A.] Americans now attach more importance to the effective use of time.
\item[B.] Americans today have more free time than earlier generations.
\item[C.] The number of hours Americans work has increased steadily.
\item[D.] More and more Americans feel pressed for time nowadays.
\end{itemize}

25.


\begin{itemize}[leftmargin=2em,itemsep=.2em,topsep=.4em]\item[A.] Our interpersonal relationships improve.
\item[B.] Our work efficiency increases greatly.
\item[C.] Our living habits are altered.
\item[D.] Our behavior is changed.
\end{itemize}

\subsection*{\textbf{Section C}}


\textbf{Directions}：\textit{In this section},\textit{ you will hear a passage three times. When the passage is read for the first time},\textit{ you should listen} \textit{carefully for its general idea. When the passage is read for the second time},\textit{ you are required to fill in the blanks with the} \textit{exact words you have just heard. Finally, when the passage is read for the third time}，\textit{you should check what you have} \textit{written.} 注意：此部分试题请在答题卡1 上作答。


The first copyright law in the United States was passed by Congress in 1790. In 1976 Congress enacted the latest copyright law, \uline{26} the technological developments that had occurred since the passage of the Copyright Act of 1909. For example, in 1909, anyone who wanted to make a single copy of a \uline{27} work for personal use had to do so by hand. The very process \uline{28} a limitation on the quantity of materials copied. Today, a photocopier can do the work in seconds: the limitation has disappeared. The 1909 law did not provide full protection for films and sound recordings, nor did it \uline{29} the need to protect radio and television. As a result, \uline{30} of the law and abuses of the intent of the law have lessened the \uline{31 }rewards of authors, artists, and producers. The 1976 Copyright Act has not prevented these abuses fully, but it has clarified the legal rights of the injured parties and given them an \uline{32} for remedy.


Since 1976 the Act has been \uline{33 }to include computer software, and guidelines have been adopted for fair use of television broadcasts. These changes have cleared up much of the confusion and conflict that followed \uline{34} the 1976 legislation.


The fine points of the law are decided by the courts and by acceptable common practice over time. As these decisions and agreements are made, we modify our behavior accordingly. For now, we need to\uline{ 35 }the law and its guidelines as accurately as we can and to act in a fair manner.


\subsection*{\textbf{Part} III \textbf{Reading Comprehension} \textbf{(40 minutes)}}


\subsection*{\textbf{Section A}}


\textbf{Directions}：\textit{In this section},\textit{ there is a passage with ten blanks. You are required to select one word for each blank from a list} \textit{of choices given in a word bank following the passage}.\textit{ Read the passage through carefully before making your choices.}


\subsection*{\textit{Each choice in the bank is identified by a letter. Please mark the corresponding letter for each item on}\textbf{\textit{ Answer Sheet 2}}}


\textit{with a single line through the centre. You may not use any of the words in the bank more than once}.


\textbf{Questions 36 to 45 are based on the following passage.}


Fear can be an effective way to change behavior. One study compared the effects of high-fear and low-fear appeals on changes in attitudes and behaviors related to dental\textit{ hygiene} (卫生）. One group of subjects was shown awful pictures of\uline{ 36 }teeth and diseased gums；another group was shown less frightening materials such as plastic teeth, charts, and graphs. Subjects who saw the frightening materials reported more anxiety and a greater\uline{ 37 }to change the way they took care of their teeth than the low-fear group did.


But were these reactions actually \uline{38} into better dental hygiene practices? To answer this important question, subjects were called back to the laboratory on two \uline{39} (five days and six weeks after the experiment). They chewed\textit{ disclosing wafers}(牙疾诊断片) that give a red stain to any uncleaned areas of the teeth and thus provided a direct \uline{40} of how well they were really taking care of their teeth. The result showed that the high-fear appeal did actually result in greater and more \uline{41 }changes in dental hygiene. That is, the subjects \uline{42} to high-fear warnings brushed their teeth more \uline{43} than did those who saw low-fear warnings.


However , to be an effective persuasive device it is very important that the message not be too frightening and that people be given \uline{44} guidelines to help them to reduce the cause of the fear. If this isn’t done, they may reduce their anxiety by denying the message or the \uline{45} of the communicator. If that happens, it is unlikely that either attitude or behavior change will occur. 注\CJKunderline{意：此部分试题请在答题卡2 上作答。}


\begin{itemize}[leftmargin=2em,itemsep=.2em,topsep=.4em]\item[A.] accustomed
\item[B.] carefully
\item[C.] cautiously
\item[D.] concrete
\item[E.] credibility
\item[F.] decayed
\item[G.] desire
\item[H.] dimensions
\item[I.] eligible
\item[J.] exposed
\item[K.] indication
\item[L.] occasions
\item[M.] permanent
\item[N.] sensitivity
\item[O.] translated
\end{itemize}

\subsection*{\textbf{Section B}}


Directions：\textit{In this section, you are going to read a passage with ten statements attached to it. Each statement contains} \textit{information given in one of the paragraphs. Identify the paragraph from which the information is derived. You may choose a} \textit{paragraph more than once. Each paragraph is marked with a letter. Answer the questions by marking the corresponding}


\subsection*{\textit{letter on}\textbf{\textit{ Answer Sheet 2.}}}


The Street-Level Solution


A) When I was growing up, one of my father’s favorite sayings (borrowed from the humorist Will Rogers) was：“It isn’t what we don’t know that causes the trouble；\textit{it ’ s what we think we know} that just isn’t so. One of the main insights to be taken from the 100 000 Homes Campaign and its strategy to end chronic homelessness is that, until recently, our society thought it understood the nature of homelessness, but it didn’t.


B) That led to a series of mistaken assumptions about why people become homeless and what they need. Many of the errors in our homelessness policies have stemmed from the conception that the homeless are a homogeneous group. It’s only in the past 15 years that organizations like Common Ground, and others, have taken a street-level view of the problem—distinguishing the “episodically homeless” from the “chronically homeless” in order to understand their needs at an individual level. This is why we can now envisage a different approach—and get better results.


C) Most readers expressed support for the effort, although a number were skeptical, and a few utterly dismissive, about the chances of long-term homeless people adapting well to housing. This is to be expected; it’s hard to imagine what we haven’t yet seen. As Niccolo Machiavelli wrote in\textit{ The Prince}, one of the major obstacles in any effort to advance systemic change is the “incredulity of men，” which is to say that people “do not readily believe in new things until they have had a long experience of them.” Most of us have witnessed homeless people on the streets for decades. Few have seen formerly homeless people after they have been housed successfully. We don’t have reference points for that story. So we generalize from what we know-or think we know.


D) But that can be misleading, even to experts. When I asked Rosanne Haggerty, founder of Common Ground, which currently operates 2 310 units of supportive housing (with 552 more under construction), what had been her biggest surprise in this work, she replied : “Fifteen years ago, I would not have believed that people who had been so broken and stuck in homelessness could thrive to the degree that they do in our buildings. ’’ And Becky Kanis, the campaign’s director, commented : ‘There is this sense in our minds that someone who’s on the streets is almost in their DNA different from someone who has a house. The campaign is creating a first-hand experience for many people that that is really not the case.


E) One of the startling realizations that I had while researching this column is that anybody could become like a homeless person —all it takes is a\textit{ traumatic} (创伤的）brain injury. A bicycle fall, a car accident, a slip on the ice, or if you’re a soldier, a head wound—and your life could become unrecognizable. James O’Connell, a doctor who has been treating the most vulnerable homeless people on the streets of Boston for 25 years, estimates that 40 percent of the long-term homeless people he’s met had such a brain injury. “For many it was a head injury prior to the time they became homeless，” he said. “They became unpredictable. They’d have mood swings, fits of explosive behavior. They couldn’t hold onto their jobs. Drinking made them feel better. They’d end up on the streets.


F) Once homeless people return to housing, they’re in a much better position to rebuild their lives. But it’s important to note that housing alone is not enough. As with many complex social problems, when you get through the initial crisis, you have another problem to solve which is no less challenging. But it is a better problem.


G) Over the past decade, O’Connell has seen this happen. spend half my time on the streets or in the hospital and the other half making house calls to people who lived for years on the streets,” he said. “So from a doctor’s point of view it’s a delightful switch, but it’s not as if putting someone in housing is the answer to addressing all of their problems. It’s the first step. ”


H) Once in housing, formerly homeless people can become isolated and lonely. If they’ve lived on the streets for years, they may have acquired a certain standing as well as a sense of pride in their survival skills. Now indoors, those aspects of their identity may be stripped away. Many also experience a profound disorientation at the outset. “If you’re homeless for more than six months, you kind of lose your bearings, ”says Haggerty. “Existence becomes not about overcoming homelessness but about finding food, begging, looking for a job to survive another day. The whole process of how you define stability gets reordered. ”


I) Many need regular, if not continuous, support with mental health problems, addictions and illnesses— and, equally important, assistance in the day-to-day challenges of life, reacquainting with family, building relationships with neighbors, finding enjoyable activities or work, managing finances, and learning how to eat healthy food.


J) For some people, the best solution is to live in a


\textit{communal}(集体)residence, with special services. This isn’t available everywhere, however. In Boston, for example, homeless people tend to be scattered in apartments throughout the city.


K) Common Ground’s large residences in New York offer insight into the possibilities for change when homeless people have a rich array of supports. In addition to more traditional social services, residents also make use of communal gardens, classes in things like cooking, yoga, theatre and photography, and job placement. Last year, 188 formerly homeless tenants in four of Common Ground\textasciicircum{} residences, found jobs.


L) Because the properties have many services and are well-managed, Haggerty has found posthousing problems to be


surprisingly rare. In the past 10 years, there have been only a handful of incidents of quarrels between tenants. There is very little\textit{ graffiti} (涂鸦）or\textit{ vandalism} (破坏）. And the turnover is almost negligible. In the Prince George Hotel in New York, which is home to 208 formerly homeless people and 208 low-income tenants, the average length of tenancy is close to seven years. (All residents pay 30 percent of their income for rent；for the formerly homeless, this comes out of their government benefits.) When people move on, it is usually because they’ve found a preferable apartment.


M) “Tenants also want to participate in shaping the public areas of the buildings，” said Haggerty. “They formed a gardening committee. They want a terrace on the roof. Those are things I didn’t count on.’’ The most common tenant demand? People always want more storage space—but that’s true of every New Yorker,” she adds. “In many ways, we’re a lot like a normal apartment building. Our tenants look like anyone else. ”


N) As I mentioned, homelessness is a catch-all for a variety of problems. A number of readers asked whether the campaign will address family homelessness, which has different causes and requires a different solution. I’ve been following some of the promising ideas emerging to address and prevent family homelessness. Later in 2011, I’ll explore these ideas in a column. For now, I’ll conclude with an update on the 100 000 Homes Campaign. Since Tuesday, New Orleans and a few other communities have reported new results. The current count of people housed is 7 043.


注意：此部分试题请在答题卡2 上作答。


46. Tenants in Common Ground’s residences all want more room for storage.


47. Homes Campaign provides first-hand proof that the homeless are not what they were once believed to be.


48. Common Ground’s residences are well-managed and by and large peaceful.


49. Housing the homeless is only the first step to solving all their problems.


50. A large percent of the chronically homeless have suffered from brain injury.


51. After being housed many homeless people become confused at first as to how to deal with life off the street.


52. Some people think the best way to help the homeless is to provide them with communal housing.


53. The homeless with health problems should be given regular support in their daily lives.


54. Until recently American society has failed to see what homelessness is all about.


55. Many formerly homeless tenants in New York’s Common Ground’s residences got hired.


\subsection*{\textbf{Section C}}


\textbf{Directions} ：\textit{There are 2 passages in this section}.\textit{ Each passage is followed by some questions or unfinished statements. For} \textit{each of them there are four choices marked A).B.)C) and D) . You should decide on the best choice and mark the}


\subsection*{\textit{corresponding letter on}\textbf{\textit{ Answer Sheet 2}}\textit{ with a single line through the centre.}}


\subsection*{\textbf{Passage One}}


\textbf{Questions 56 to 60 are based on the following passage.}


Technology can make us smarter or stupider, and we need to develop a set of principles to guide our everyday behavior and make sure that tech is improving and not hindering our mental processes. One of the big questions being debated today is: What kind of information do we need to have stored in our heads, and what kind can we leave “in the cloud，” to be accessed as necessary?


An increasingly powerful group within education are championing “digital literacy”. In their view, skills beat knowledge, developing “digital literacy” is more important than learning mere content, and all facts are now Google-able and therefore unworthy of committing to memory. But even the most sophisticated digital literacy skills won’t help students and workers navigate the world if they don’t have a broad base of knowledge about how the world actually operates. If you focus on the delivery mechanism and not the content, you’re doing kids a disservice.


Indeed, evidence from cognitive science challenges the notion that skills can exist independent of factual knowledge. Data from the last thirty years leads to a conclusion that is not scientifically challengeable：thinking well requires knowing facts, and that’s true not only because you need something to think about. The very processes that teachers care about most —critical thinking processes—are intimately\textit{ intertwined}(交织）with factual knowledge that is stored in long-term memory.


In other words, just because you can Google the date of Black Tuesday doesn’t mean you understand why the Great Depression happened or how it compares to our recent economic slump. There is no doubt that the students of today, and the workers of tomorrow, will need to innovate, collaborate and evaluate. But such skills can’t be separated from the knowledge that gives rise to them. To innovate, you have to know what came before. To collaborate, you have to contribute knowledge to the joint venture. And to evaluate, you have to compare new information against knowledge you’ve already mastered.


So here’s a principle for thinking in a digital world, in two parts. First, acquire a base of factual knowledge in any domain in which you want to perform well. This base supplies the essential foundation for building skills, and it can’t be\textit{ outsourced} (外包） to a search engine.


Second, take advantage of computers’ invariable memory, but also the brain’s elaborative memory. Computers are great when you want to store information that shouldn’t change. But brains are the superior choice when you want information to change, in interesting and useful ways：to connect up with other facts and ideas, to acquire successive layers of meaning, to steep for a while in your accumulated knowledge and experience and so produce a richer mental brew. 注意：此部分试题请在答题卡2 上作答。


56. What is the author’s concern about the use of technology?


\begin{itemize}[leftmargin=2em,itemsep=.2em,topsep=.4em]\item[A.] It may leave knowledge “in the cloud”.
\item[B.] It may misguide our everyday behavior.
\item[C.] It may cause a divide in the circles of education.
\item[D.] It may hinder the development of thinking skills.
\end{itemize}

57. What is the view of educators who advocate digital literacy?


\begin{itemize}[leftmargin=2em,itemsep=.2em,topsep=.4em]\item[A.] It helps kids to navigate the virtual world at will.
\item[B.] It helps kids to broaden their scope of knowledge.
\item[C.] It increases kids’ efficiency of acquiring knowledge.
\item[D.] It liberates kids from the burden of memorizing facts.
\end{itemize}

58. What does evidence from cognitive science show?


\begin{itemize}[leftmargin=2em,itemsep=.2em,topsep=.4em]\item[A.] Knowledge is better kept in long-term memory.
\item[B.] Critical thinking is based on factual knowledge.
\item[C.] Study skills are essential to knowledge acquisition.
\item[D.] Critical thinking means challenging existing facts.
\end{itemize}

59. What does the author think is key to making evaluations?


\begin{itemize}[leftmargin=2em,itemsep=.2em,topsep=.4em]\item[A.] Gathering enough evidence before drawing conclusions.
\item[B.] Mastering the basic rules and principles for evaluation.
\item[C.] Connecting new information with one’s accumulated knowledge.
\item[D.] Understanding both what has happened and why it has happened.
\end{itemize}

60. What is the author’s purpose in writing the passage?


\begin{itemize}[leftmargin=2em,itemsep=.2em,topsep=.4em]\item[A.] To warn against learning through memorizing facts.
\item[B.] To promote educational reform in the information age.
\item[C.] To explain human brains5 function in storing information.
\item[D.] To challenge the prevailing overemphasis on digital literacy.
\end{itemize}

\subsection*{\textbf{Passage Two}}


\textbf{Questions 61 to 65 are based on the following passage.}


America’s recent history has been a persistent tilt to the West—of people, ideas, commerce and even political power. California and Texas are the twin poles of the West, but very different ones. For most of the 20th century the home of Silicon Valley and Hollywood has been the brainier and trendier of the two. Texas has trailed behind：its stereotype has been a conservative Christian in cowboy boots. But twins can change places. Is that happening now?


It is easy to find evidence that California is in a panic. At the start of this month the once golden state started paying creditors in\textit{ IOUs} (欠条).The gap between projected outgoings and income for the current\textit{ fiscal} (财政的）year has leapt to a horrible \$26 billion. With no sign of a new budget to close this gulf, one credit agency has already downgraded California’s debt. As budgets are cut, universities will let in fewer students, prisoners will be released early and schemes to protect the vulnerable will be rolled back.


By contrast, Texas has coped well with the recession, with an unemployment rate two points below the national average and one of the lowest rates of housing repossession. In part this is because Texan banks, hard hit in the last property bust, did not overexpand this time. Texas also clearly offers a different model, based on small government. It has no state capital-gains or income tax, and a business- friendly and immigrant-tolerant attitude. It is home to more\textit{ Fortune} 500 companies than any other state.


Despite all this, it still seems too early to hand over America’s future to Texas. To begin with, that lean Texan model has its own problems. It has not invested enough in education, and many experts rightly worry about a “lost generation” of mostly Hispanic Texans with insufficient skills for the demands of the knowledge economy.


Second, it has never paid to bet against a state with as many inventive people as California. Even if Hollywood has gone into depression, it still boasts an unequalled array of sunrise industries and the most brisk venture-capital industry on the planet. The state also has an awesome ability to reinvent itself —as it did when its defence industry collapsed at the end of the cold war.


The truth is that both states could learn from each other. Texas still lacks California’s great universities and lags in terms of culture. California could adopt not just Texas’s leaner state, but also its more\textit{ bipartisan} (两党的）approach to politics. There is no perfect model of government: it is America’s genius to have 50 public-policy laboratories competing to find out what works best. 注意：此部分试题请在答题卡2 上作答。


61. What does the author say about California and Texas in Paragraph 1?


\begin{itemize}[leftmargin=2em,itemsep=.2em,topsep=.4em]\item[A.] They have been competing for the leading position.
\item[B.] California has been superior to Texas in many ways.
\item[C.] They are both models of development for other states.
\item[D.] Texas’s cowboy culture is less known than California’s.
\end{itemize}

62. What does the author say about today’s California?


\begin{itemize}[leftmargin=2em,itemsep=.2em,topsep=.4em]\item[A.] Its debts are pushing it into bankruptcy.
\item[B.] Its budgets have been cut by \$26 billion.
\item[C.] It is faced with a serious financial crisis.
\item[D.] It is trying hard to protect the vulnerable.
\end{itemize}

63. In what way is Texas different from California?


\begin{itemize}[leftmargin=2em,itemsep=.2em,topsep=.4em]\item[A.] It practices small government.
\item[B.] It is home to traditional industries.
\item[C.] It has a large Hispanic population.
\item[D.] It has an enviable welfare system.
\end{itemize}

64. What problem is Texas confronted with?


\begin{itemize}[leftmargin=2em,itemsep=.2em,topsep=.4em]\item[A.] Its Hispanic population is mostly illiterate.
\item[B.] Its sunrise industries are shrinking rapidly.
\item[C.] Its education cannot meet the needs of the knowledge economy.
\item[D.] Its immigrants have a hard time adapting to its cowboy culture.
\end{itemize}

65. What do we learn about American politics from the passage?


\begin{itemize}[leftmargin=2em,itemsep=.2em,topsep=.4em]\item[A.] Each state has its own way of governing.
\item[B.] Most states favor a bipartisan approach.
\end{itemize}

\begin{itemize}[leftmargin=2em,itemsep=.2em,topsep=.4em]\item[C.] Parties collaborate in drawing public policies.
\item[D.] All states believe in government for the people.
\end{itemize}

\subsection*{\textbf{Part} IV \textbf{Translation} \textbf{(30 minutes)}}


\textbf{Directions}：\textit{For this part},\textit{ you are allowed 30 minutes to translate a passage from Chinese into English}.\textit{ You should write}


\subsection*{\textit{your answer on}\textbf{\textit{ Answer Sheet 2.•}}}


中文热词通常反映社会变化和文化，有些在外国媒体上愈来愈流行。例如,土豪（tuhao)和大妈（dama)都是老词， 但已获取了新的意义。


土豪以前指欺压佃户和仆人的乡村地主，现在用于指花钱如流水或喜欢炫耀财富的人。也就是说，土豪有钱，但没有品位。大妈是对中年妇女的称呼，但现在特指不久前金价下跌时大量购买黄金的中国妇女。


土豪和大妈可能会被收入新版牛津(Oxford)英语词典。至今已有约120 个中文词被加进了牛津英语词典,成了英语语言的一部分。 注意：此部分试题请在答题卡2 上作答。

\end{document}